\section{Feature Module Details}\label{app:method}

\textbf{DCT double quantisation.} Under protocols B and R the last compression is known (Q85), so the module uses
the exact IJG Q85 quantisation table; a blind estimate of the last quantisation step is used only when it is
unknown (protocol A, or arbitrary input files). Periodicity is tested on the log-histogram of each frequency's
coefficients after removing a quadratic trend. The residual is projected on every candidate period from 2 to half
the histogram range, including fractional periods. The largest normalised power is compared with its
\emph{permutation null}: the same maximum computed after shuffling the residual across bins, 200 times per
histogram. This gives a calibrated p-value whatever the histogram's width or tails, and however many periods are
tried. Periodicity is declared at $p\leq\alpha$, with $\alpha$ tuned in Section~\ref{sec:params}. On dedicated test
cases the false-positive rate at $\alpha = 0.01$ is 0--3\% (0.4\% on real history-free images), against 16--48\%
for an earlier fixed threshold. Genuinely double-compressed test images are flagged at 75\% of frequencies, but
overall detection on simulated double compression is only about 56\%: some regular patterns are missed, e.g.\
$q_1/q_2 = 10/4$ (a period of 2.5 bins) was detected in 0/30 simulations, while 9/4, 7/3 and 8/3 were detected in
30/30. The per-coefficient posterior compares the observed histogram with a uniform one inside a window spanning a
whole number of periods.

\textbf{Blind quantisation-step estimator.} The estimator picks the largest step that is a local maximum of the
integer-fit score and reaches 80\% of what pixel-rounding noise allows ($\sigma = 0.5$, relaxed to 1.0 for
saturated images). When the first step is an exact multiple of the last (e.g.\ 9 and 3), double compression cannot
be detected blindly at that frequency.

\textbf{Mirrored copies.} Several CASIA copy-moves are horizontally flipped duplicates, which neither SIFT
descriptors nor block DCT features match directly. Both detectors therefore also match against the mirrored image.
For keypoints, descriptors are recomputed at the mirrored keypoint positions. For blocks, the shift measured in
(image, mirror) coordinates is constant for a mirrored copy, so ordinary shift voting still applies; the mirrored
image's block grid is offset so that it aligns with the image's block grid.

\textbf{Pair orientation.} A translation has a direction, so direct keypoint pairs are oriented left-to-right
before RANSAC. Unoriented pairs split one true copy into two opposite transforms, and RANSAC then rejected genuine
copies. Mirrored pairs are oriented left-to-right as well: a mirrored copy with a vertical offset is a glide
reflection, whose inverse is a different transform, so the same split would occur. One RANSAC fit is run per
matching mode (direct and mirrored), so the keypoint detector recovers at most one copied region per mode.

\textbf{Verification.} The synthetic forgeries with a known region are: for ELA and the histogram module, an
uncompressed region inside a Q60 image; for noise, a noisier region; for JPEG ghosts, a region previously saved at
Q60 inside a Q90 image, re-saved at Q95; for DCT, a single-compressed region inside a double-compressed image; for
edges, an unblurred region in a blurred image of the same content; and for copy-move, plain and mirrored copies.
Contract tests cover tiny ($8\times8$), narrow ($200\times13$), flat and black images, and check determinism. The
modules are implemented in \texttt{src/forgery/features/} and tested in \texttt{tests/test\_features.py}; the
synthetic tests run both with the default parameters and with the tuned parameters from \texttt{config.yaml}.
